% \iffalse meta-comment
%
% hit-thesis-abstract.dtx 是中英文摘要页
%
% \fi
%
% \section{中英文摘要页}
%
%    \begin{macrocode}
%<*thesis-abstract>
%    \end{macrocode}
%
% 生成参考文献和关键字。
%    \begin{macrocode}
\box_new:N \hit@keywords@box
\NewDocumentCommand \hit@put@keywords { m m }
  {
    \group_begin:
      \setbox \hit@keywords@box = \hbox {#1}
      \noindent \hangindent \wd \hit@keywords@box \hangafter 1
      \box \hit@keywords@box #2 \par
    \group_end:
  }

%    \end{macrocode}
%
% \changes{v3.0.0}{2020/05/21}{英文版目录之前有page行}
% \changes{v3.0.0}{2020/05/21}{博后关键字加粗}
% \changes{v3.2a}{2026/08/07}{\cs{hit@make@abstract} 迁到 expl3。写进目录的
%   \texttt{\string~} 换成 \cs{nobreakspace}\{\}，因为 expl3 里 \texttt{\string~}
%   是普通空格}
%    \begin{macrocode}
\cs_new_protected:Npn \hit@make@abstract
  {
    \hit@clear@page
    \bool_if:NT \g__hit_en_bool
      {
        \addtocontents { toc } { \nobreakspace { } \hfill \textbf { \hit@contents@page@heading } \par }
        \addtocontents { \ext@table }
          { { \noindent \sihao \textbf { \hit@contents@table@heading } \nobreakspace { } \hfill \textbf { \hit@contents@page@heading } } \par }
        \addtocontents { \ext@figure }
          { { \noindent \sihao \textbf { \hit@contents@figure@heading } \nobreakspace { } \hfill \textbf { \hit@contents@page@heading } } \par }
      }
    \hit@appendix@chapter * [ \cabstractcname ]
      { \hit@spread@by:nn { \hit@heading@spread } { \cabstractcname } }
      [ \cabstractename ]
    \pagestyle { hit@headings }
    \pagenumbering { Roman }
    \hit@abstract@zh
    % 正文与关键词之间按范例是一个空段(单独回车一行),\hitenter 模拟之
    % (1-8 页版 p4,36.8 = 19.4 + 17.25)
    \hitenter
    \hit@put@keywords
      { \bool_if:NT \g__hit_postdoc_bool { \bfseries } \heiti \hit@keywords@field@label@zh }
      { \hit@keywords@zh }
    \hit@clear@page
    \hit@appendix@chapter * { \eabstractcname } [ \eabstractename ]
    % 英文摘要正文按范例排:Word 的 1.25 倍自动行距按行内字体自然行高算,
    % 纯西文段 TNR 12bp 实测行距 17.25;字符网格的行末余量 \rightskip 只属于
    % 汉字网格行,西文段排满真版心(例右缘 510.4);汉字行撑竿会把行距顶回
    % 网格,组内停用;范例导出勾了不断字,断词罚分同关
    \group_begin:
      % 订阅拉丁段制度(行距/版宽/撑竿/断词,见 hit-format)
      \hit@format@latin@par:
      \hit@abstract@en
      \hitenter
      \hit@put@keywords { \textbf { \hit@keywords@field@label@en \enskip } } { \hit@keywords@en }
    \group_end:
  }
%    \end{macrocode}
%
% \changes{v3.2a}{2026/08/10}{\env{denotation} 收成 \cs{item} 列表，
%   原先自己写表格的写法移到 \env{denotation*}}
% 符号表。
%
% 原先这个环境只管标题和表号，正文要用户自己摆一张表，列宽、字号、跨页都得手调。
% 现在环境本身就是一个 \pkg{enumitem} 的 \texttt{description} 列表，一行一个符号：
% \cs{item}\oarg{符号}\meta{说明}。第一列宽度默认 2.5cm，
% \env{denotation} 的可选参数改它。
%
% 规范里允许把符号分成几张表、各有各的列，列表排不出来，所以保留
% \env{denotation*}：标题和表号照旧，正文由用户自己写，跟 v3.1e 一模一样。
%
% 环境本身认不出用户想要哪一种，所以在开头探一下：正文第一个记号是 \cs{item} 就开
% 列表，不是就当旧写法，照常编译、提示加个星号。前面的空格和空行先吃掉，不然
% \cs{begin}\marg{denotation} 后面空一行就会被误判成旧写法。
%    \begin{macrocode}
\RequirePackage{enumitem}
\newlist { hit@denotation } { description } { 1 }
\setlist [ hit@denotation ]
  {
    nosep ,
    font = \normalfont ,
    align = left ,
    leftmargin = ! ,% 下面三个长度之和
    labelindent = 0pt ,
    labelwidth = 2.5cm ,
    labelsep* = 0.5cm ,
    itemindent = 0pt ,
  }
\bool_new:N \g__hit_denotation_list_bool
\tl_new:N \g__hit_denotation_width_tl
\cs_new_protected:Npn \__hit_denotation_open:
  {
    \hit@clear@page
    \hit@appendix@chapter * { \hit@denotation@heading@zh } [ \hit@denotation@heading@en ]
    \setcounter { table } { 0 }
    \cs_set:Npn \thetable { \arabic { table } }% 表编号为 1
  }
\cs_new_protected:Npn \__hit_denotation_close:
  {
    \cs_set:Npn \thetable { \arabic { chapter } - \arabic { table } }% 表编号为 7-1
    \setcounter { table } { 0 }
  }
\cs_new_protected:Npn \__hit_denotation_peek:
  {
    \peek_meaning_remove_ignore_spaces:NTF \par
      { \__hit_denotation_peek: }
      {
        \peek_meaning_ignore_spaces:NTF \item
          {
            \bool_gset_true:N \g__hit_denotation_list_bool
            \begin { hit@denotation } [ labelwidth = \g__hit_denotation_width_tl ]
          }
          {
            \bool_gset_false:N \g__hit_denotation_list_bool
            \msg_warning:nn { hithesis } { denotation-no-item }
          }
      }
  }
\NewDocumentEnvironment { denotation } { O{2.5cm} }
  {
    \tl_gset:Nn \g__hit_denotation_width_tl {#1}
    \__hit_denotation_open:
    \__hit_denotation_peek:
  }
  {
    \bool_if:NT \g__hit_denotation_list_bool { \end { hit@denotation } }
    \__hit_denotation_close:
  }
%    \end{macrocode}
%
% 带星号的不开列表，正文完全交给用户。可选参数收下不用，这样 v3.1e 里
% \cs{begin}\marg{denotation}\oarg{2.5cm} 的写法加个星号就能直接跑。
%    \begin{macrocode}
\NewDocumentEnvironment { denotation* } { O{2.5cm} }
  { \__hit_denotation_open: }
  { \__hit_denotation_close: }
%    \end{macrocode}
%
%    \begin{macrocode}
%</thesis-abstract>
%    \end{macrocode}
%
% \Finale
